% Second paper: the submask context family, the collapse of the 19-round cost,
% and the computed 20-round preimages.  Companion to paper_r19_final.tex
% (the base attack) and paper_r20_barrier.tex (the historical barrier note).
% This paper supersedes their cost figures where stated.
%
% Every number in this paper is traceable to a file under
% experiments/submask_2026-09-03/ in the repository; see Appendix B.
%
\documentclass[11pt,a4paper]{article}

\usepackage{amsmath,amssymb,amsthm}
\usepackage{booktabs}
\usepackage{xcolor}
\usepackage[hidelinks]{hyperref}
\Urlmuskip=0mu plus 1mu\relax
\hypersetup{pdftitle={Context Shaping for 19- and 20-Round SHA-256 Compression Preimages},
            pdfauthor={Kameldip Singh Basra},
            pdfkeywords={SHA-256, preimage attack, reduced-round, compression function, cryptanalysis}}
\usepackage{geometry}
\geometry{margin=1in}
\usepackage[protrusion=true,expansion=false]{microtype}
\usepackage{array}

\newtheorem{theorem}{Theorem}
\newtheorem{corollary}{Corollary}
\newtheorem{definition}{Definition}
\newtheorem{lemma}{Lemma}
\newtheorem{proposition}{Proposition}
\newtheorem{remark}{Remark}

\newcommand{\Sig}[1]{\Sigma_{#1}}
\newcommand{\sig}[1]{\sigma_{#1}}
\newcommand{\Maj}{\mathrm{Maj}}
\newcommand{\Ch}{\mathrm{Ch}}
\newcommand{\ROTR}{\mathrm{ROTR}}
\newcommand{\SHR}{\mathrm{SHR}}
\newcommand{\mmod}{\bmod\,2^{32}}
\newcommand{\hex}[1]{\texttt{#1}}

\begin{document}

\title{\textbf{Context Shaping for 19- and 20-Round\\
       SHA-256 Compression Preimages}}

\author{Kameldip Singh Basra \\[4pt]
        \small Independent Researcher \\
        \small \texttt{kameldipbasra@gmail.com} \\
        \small ORCID: \href{https://orcid.org/0009-0001-0977-4614}{0009-0001-0977-4614}}

\date{September 2026}

\maketitle

\begin{abstract}
\noindent
We study one-block SHA-256 compression reduced to 19 or 20 rounds, with the
standard initial value, feed-forward retained, all 256 output bits fixed, and
an unrestricted 512-bit input block.  A global table for
$u\mapsto\sigma_0(u)-u$ underlies an earlier 19-round construction.  We show
that choosing internal context words so that $a_4=a_5$ and
$e_8=e_9=\hex{0xFFFFFFFF}$ makes its three schedule constraints triangular:
three table lookups replace a fixed-point iteration.  The same conditions
reduce the remaining consistency test to
$\Maj(a_4,a_3,a_2)=a_3$, which holds for exactly $3^{32}$ of the $2^{64}$
pairs $(a_2,a_3)$.  Under an explicit lookup-distribution model this predicts
one 19-round preimage per 39{,}124 swept values; three recorded runs give one
per 38{,}878.  A 64-target run produced 13{,}840 forward-verified preimages in
151.7 seconds.  The often-quoted 12 ms figure is a marginal estimate obtained
from measured single-core throughput and yield after a reusable 16-GiB table
has been built, not a cold-start latency.

At 20 rounds, candidates must also satisfy an exact fourth schedule
constraint.  Existing residual tests are consistent with the working model
that this condition contributes a factor $2^{-32}$; they do not prove
uniformity at exact zero or rule out all alternative policies.  With three
stored roots per table value, the model predicts a mean cost of about
$2^{45.4}$ swept values, or about 15 hours at the measured A100 production
rate.  We report two computed, forward-verified examples: one
for a solver-generated target after 2.45 GPU-hours and one for the fixed
all-ones target after 16.16 GPU-hours.  The six recorded production runs total
about 71.7 GPU-hours.  The method combines known choice- and
majority-function degeneracies.  Our negative 21-round observations apply
only to the tested construction and do not establish a lower bound.  These
results concern reduced-round compression, not padded or full SHA-256.
\end{abstract}

\section{Introduction}
\label{sec:intro}

A \emph{preimage} of the SHA-256 compression function reduced to $R$ rounds is
a 512-bit block $W_0,\dots,W_{15}$ such that the $R$-round compression of the
standard initial value with that block, plus the feed-forward, equals a given
256-bit digest.  We retain feed-forward but impose no SHA-256 padding format on
the block.  Thus the object is the one-block reduced-round compression
function, not the standard padded hash function.  The fixed initial value also
distinguishes it from a free-start or pseudo-preimage problem.  All target bits
are fixed; this is not a partial-output or collision problem.  This is the
model used by Zaikin~\cite{Zaikin2026} and by our earlier
paper~\cite{Basra2026R19}.

Two lines of work address this problem.  The meet-in-the-middle line reaches
high step counts at costs just below brute force: 43 steps at $2^{254.9}$
\cite{AokiGuo2009}, 45 at $2^{255.5}$ \cite{KRS2012}, and, since February
2026, 47 at $2^{255.6}$ \cite{GuoLiLiuWangZhang2026}.  None of those attacks
has been executed in full (the largest computed fragment is discussed in
Section~\ref{sec:priority}), and the cheapest figure anywhere in that line is $2^{240}$
at 24 steps~\cite{IsobeShibutani2009}.  The practical line, which actually
computes preimages, has moved slowly: 16 rounds by SAT solving in
2012~\cite{Homsirikamol2012}, 17 and 18 rounds for the all-zero digest in
2025~\cite{Davydov2025}, and 19 rounds in January 2026, when
Zaikin~\cite{Zaikin2026} inverted the all-ones digest in 18~h~33~min on 192
CPU cores and wrote that ``it is likely that additional effective algorithmic
techniques are required to invert \ldots\ 20-round SHA-256 compression
functions'' (\cite{Zaikin2026}, \S8).  Our earlier paper~\cite{Basra2026R19} reached 19 rounds by a
different, algebraic route.  Its screened/control campaign gave a pooled
point estimate of about two GPU-hours per success, with substantial sampling
uncertainty and no events in its random-context control arm.  A companion
note~\cite{Basra2026R20} modelled the 20-round extension as a further
$2^{-32}$ schedule filter.

This paper supersedes the performance estimates of our earlier manuscripts
for this chosen-context family and reports two computed and verified preimages
of the 20-round SHA-256 compression function.  To our knowledge, after the
dated and scoped literature search described in Section~\ref{sec:priority}, no
published work had exhibited such a full-output witness beyond 19 rounds.

\paragraph{What the base attack does.}
The attack of~\cite{Basra2026R19} fixes the last eight state words from the
digest by a backward chain, chooses the seven state words $a_4,\dots,a_{10}$
at random (the \emph{context}), sweeps $a_0$ over all $2^{32}$ values, and
recovers $a_1,a_2,a_3$ from the three message-schedule constraints on
$W_{16},W_{17},W_{18}$ using a 16-GiB table that inverts $u\mapsto\sig0(u)-u$.
Two things make it expensive.  The constraints on $W_{17}$ and $W_{18}$ are
coupled through the round function, so $a_2$ and $a_3$ are found by a
fixed-point iteration that converges for roughly one $C_0$ survivor in
$2^{16}$ ($1.3$--$1.9\times10^{-5}$ per survivor in~\cite{Basra2026R19}).  And the constraint on $W_{16}$ needs $W_9$, which depends on the
still-unknown $a_2,a_3$; the attack freezes them at zero, solves, and
afterwards tests that the true $W_9$ equals the provisional one.  The archived
experiments measure substantial, context-dependent enrichment in this final
condition; the earlier paper reports the resulting campaign rates rather than
treating the residual as uniform.

\paragraph{What this paper does.}
We choose the context instead of drawing it at random.  Three conditions, $a_4=a_5$ together with $e_8=e_9=\hex{0xFFFFFFFF}$ (the
last two each imposed by solving one round equation for one context word),
have two consequences
(Section~\ref{sec:family}).  The coupling that forced the iteration
disappears: the constraint on $W_{17}$ loses its dependence on $a_3$, so
$a_1,a_2,a_3$ are obtained by three table lookups.  And the consistency
residual collapses to $\Maj(a_4,a_3,a_2)-a_3$, a difference of two bitwise
functions whose equality is a per-bit condition
$\Maj(a_4,a_3,a_2)=a_3$, satisfied by exactly $3^{32}$ of the $2^{64}$ pairs.

\paragraph{Results.}
\begin{itemize}
\item At 19 rounds the cost per preimage is about $2^{15.25}$ swept $a_0$,
  pooled over $1.14\times10^{9}$ swept values and $29{,}398$ verified
  preimages in three runs.  The historical random-context campaign reported
  $2^{38.3}$ swept values per preimage when its screened and control arms were
  pooled, but that mixed-arm number is not an estimate for uniformly sampled
  contexts; we retain it only as a description of the earlier implementation
  (Section~\ref{sec:r19}).  The
  model-derived value is $1/(0.633673^3\cdot(3/4)^{32}) = 39{,}124$;
  the three-run pooled value is $38{,}878$ ($1{,}142{,}947{,}840$ swept $a_0$ over $29{,}398$ preimages), and the largest separate replication (Table~\ref{tab:r19}, 24 contexts) alone gives $39{,}123$.  One CPU core produces a preimage
  at a marginal modelled time of about 12~ms, derived from measured throughput
  and yield after table construction.  The all-ones digest of~\cite{Zaikin2026}
  and the all-zero benchmark of~\cite{Davydov2025} each yield about a hundred
  preimages per second.
\item At 20 rounds the fourth schedule constraint is an exact 32-bit filter in which we find no bias that a table policy or a
  context choice could exploit in the tested distributions and statistics
  (Section~\ref{sec:r20}).  Assuming its exact-zero rate is $2^{-32}$, three-root
  tables and a fused GPU kernel give a modelled mean cost of $2^{45.4}$ swept
  $a_0$, about 15~hours at the measured A100 production rate.  We computed two 20-round preimages: one for a
  digest generated by the solver, after 2.45~h on one A100, and one for the
  all-ones digest, after 16.2~h on one of two A100s searching that digest.  To our knowledge, and
  after the search described in Section~\ref{sec:priority} (as of September
  2026), no computed preimage of the SHA-256 compression function beyond 19
  rounds had previously been published.
\item At 21 rounds two additional exact 32-bit constraints remain.  Under the
  explicitly stated model that their residuals are uniform and independent,
  they contribute a factor $2^{64}$ relative to 19 rounds.  Measurements on
  214{,}583 candidates and searches over a finite menu of context families are
  consistent with that model but do not prove it or exclude other
  representations (Section~\ref{sec:barrier}).
\end{itemize}

\paragraph{What is and is not new.}
The two levers are known.  Saturating the choice function,
$\Ch(\hex{0xFFFFFFFF},y,z)=y$, is the ``absorption property'' used for message
stealing by Aoki et~al.~\cite{AokiGuo2009}; the majority degeneracy
$\Maj(x,x,y)=x$ is standard and appears in Guo
et~al.~\cite{GuoLiLiuWangZhang2026} as an auxiliary condition that absorbs
propagation through $\Maj$ (their Proposition~11 and Remark~3).  Our own companion
note~\cite{Basra2026R20} applies both levers to a different constraint and
dismisses the result because one of its conditions tied a context word to an
unknown, which would force the constants to be rebuilt per candidate.  What is
new here is the specific pairing $a_4=a_5$, which ties two \emph{context}
words together and costs nothing per candidate, its application to the
constraint that carries the iteration, and the observation that the same
choice turns a 32-bit modular consistency check into a bitwise one.  We have
found no precedent for that last effect.  The barrier argument
of~\cite{Basra2026R20} identified the remaining residual and modelled its
cost as a factor $2^{32}$ relative to the 19-round generator; what it did not
anticipate is that the 19-round generator itself could become about $2^{23}$
cheaper.
Section~\ref{sec:relation} reconciles the two manuscripts.  The individual
choice- and majority-function identities are prior art; the claimed
contribution is their particular joint use here and the resulting residual
factorisation.

\paragraph{Scope.}
Nineteen and twenty of SHA-256's 64 rounds.  Nothing in this paper bears on
full SHA-256; on full SHA-224, which shares the compression function but uses
a different initial value and truncates its output; on SHA-512, for which the
table would need $2^{64}$ entries; or on any deployed system.  The two
20-round targets are respectively image-derived and fixed all-ones.  Their
existence does not prove that every 256-bit target is attainable or that the
modelled mean applies uniformly to every target.

\section{Preliminaries}
\label{sec:prelim}

\paragraph{Notation.}
All words are 32-bit and arithmetic is modulo $2^{32}$ unless stated.
$\ROTR^n(x)$ denotes a right rotation and $\SHR^n(x)$ a logical right shift.
The SHA-256 functions of FIPS 180-4~\cite{FIPS180} are
\begin{align*}
\Sig0(x)&=\ROTR^2(x)\oplus\ROTR^{13}(x)\oplus\ROTR^{22}(x),\\
\Sig1(x)&=\ROTR^6(x)\oplus\ROTR^{11}(x)\oplus\ROTR^{25}(x),\\
\sig0(x)&=\ROTR^7(x)\oplus\ROTR^{18}(x)\oplus\SHR^3(x),\\
\sig1(x)&=\ROTR^{17}(x)\oplus\ROTR^{19}(x)\oplus\SHR^{10}(x).
\end{align*}
Also, $\Ch(e,f,g)=(e\wedge f)\oplus(\neg e\wedge g)$ and
$\Maj(a,b,c)=(a\wedge b)\oplus(a\wedge c)\oplus(b\wedge c)$.  Following
\cite{Basra2026R19}, the state after round $r$ is written $a_r$ and $e_r$,
with $a_{-1},\dots,a_{-4}$ and $e_{-1},\dots,e_{-4}$ the eight words of the
standard initial value, in order
\[
\hex{6a09e667},\ \hex{bb67ae85},\ \hex{3c6ef372},\ \hex{a54ff53a},\
\ \hex{510e527f},\ \hex{9b05688c},\ \hex{1f83d9ab},\ \hex{5be0cd19}.
\]
Round $r$ consumes $W_r$ and the standard constant $K_r$:
\begin{align*}
T_1^{(r)} &= e_{r-4}+\Sig1(e_{r-1})+\Ch(e_{r-1},e_{r-2},e_{r-3})+K_r+W_r,\\
T_2^{(r)} &= \Sig0(a_{r-1})+\Maj(a_{r-1},a_{r-2},a_{r-3}),
\end{align*}
$a_r=T_1^{(r)}+T_2^{(r)}$ and $e_r=a_{r-4}+T_1^{(r)}$.  Hence
$e_r=a_{r-4}+a_r-T_2^{(r)}$ is determined by the $a$-words alone, and every
message word is recovered from the state:
\begin{equation}
W_r = a_r - T_2^{(r)} - e_{r-4} - \Sig1(e_{r-1}) - \Ch(e_{r-1},e_{r-2},e_{r-3}) - K_r .
\label{eq:recW}
\end{equation}
The digest after $R$ rounds is the initial value plus
$(a_{R-1},a_{R-2},a_{R-3},a_{R-4},e_{R-1},\dots,e_{R-4})$.

For $t\ge16$, the message schedule is
\begin{equation}
W_t=\sig1(W_{t-2})+W_{t-7}+\sig0(W_{t-15})+W_{t-16}.
\label{eq:schedule}
\end{equation}

\paragraph{Backward chain from the target.}
Subtracting the standard initial value word by word from the target digest
determines $(a_{R-1},\ldots,a_{R-4},e_{R-1},\ldots,e_{R-4})$.  For
$r=R-1,R-2,R-3,R-4$, in that order, compute
\begin{equation}
T_2^{(r)}=\Sig0(a_{r-1})+\Maj(a_{r-1},a_{r-2},a_{r-3}),\quad
T_1^{(r)}=a_r-T_2^{(r)},\quad
a_{r-4}=e_r-T_1^{(r)}.
\label{eq:backward}
\end{equation}
Each step uses only words already known from the digest or an earlier step,
so the target fixes eight consecutive $a$-words.  This is an algebraic
inversion of four state updates; it neither chooses nor uses a target
preimage.

\paragraph{The base attack.}
Given the digest, the backward chain of~\cite{Basra2026R19} recovers
$a_{R-8},\dots,a_{R-1}$.  At $R=19$ the free words are $a_0,\dots,a_{10}$; the
attack fixes $a_4,\dots,a_{10}$ (the context), sweeps $a_0$, and must solve for
$a_1,a_2,a_3$ against the schedule constraints
\begin{equation}
C_j:\quad W_{16+j} = \sig1(W_{14+j}) + W_{9+j} + \sig0(W_{1+j}) + W_j,\qquad j=0,1,2,
\label{eq:Cj}
\end{equation}
in which every $W$ is expressed through~\eqref{eq:recW}.  Because $a_{j+1}$
enters $W_{j+1}$ with coefficient $+1$ and $W_{9+j}$ with coefficient $-1$,
$C_j$ can be written as $\sig0(W_{j+1})-W_{j+1}=\text{(terms free of
}a_{j+1})$, and a global table inverting $u\mapsto\sig0(u)-u$ yields
$W_{j+1}$, hence $a_{j+1}$, in one lookup.  The map $u\mapsto\sig0(u)-u$ is
not a bijection: its image has exactly $2{,}721{,}603{,}628$ of the $2^{32}$
values, a fraction $0.633673$, and the published table stores one root per
image value, the largest.

The difficulty is that the right-hand side of $C_1$ depends on $a_3$ and that
of $C_2$ on $a_2$, so the attack iterates $C_1\to a_2$, $C_2\to a_3$ to a fixed
point; and the right-hand side of $C_0$ contains $W_9$, which depends on
$a_2,a_3$ through
\begin{equation}
W_9 = \underbrace{T_1^{(9)}-K_9}_{\text{context}}
      - e_5 - \Sig1(e_8) - \Ch(e_8,e_7,e_6),\qquad
e_5 = a_1 + a_5 - \Sig0(a_4) - \Maj(a_4,a_3,a_2),
\label{eq:W9}
\end{equation}
with $e_6=a_2+T_1^{(6)}$ and $e_7=a_3+T_1^{(7)}$.  The attack evaluates $W_9$
at $a_2=a_3=0$, a provisional value $\hat W_9$, solves $C_0$ with it, and
finally tests
\begin{equation}
\varepsilon \;=\; W_9^{\mathrm{conv}}(a_2,a_3) - W_9^{\mathrm{conv}}(0,0)
\;\stackrel{?}{=}\; 0,
\label{eq:eps}
\end{equation}
where $W_9^{\mathrm{conv}}$ is the $a_1$-free part of~\eqref{eq:W9}.  The
measured rates in~\cite{Basra2026R19} are about $1.6\times10^{-5}$ fixed
points per $C_0$ survivor and one success of~\eqref{eq:eps} per $2^{22}$
fixed points, giving about one preimage per 80 attempts of $2^{32}$ swept
$a_0$ each, i.e.\ $2^{38.3}$ swept $a_0$ per preimage.\footnote{The pooled campaign estimate of~\cite{Basra2026R19},
$\hat\mu=3/240$ per attempt, mixes a screened arm with a control arm that
produced no events, and carries a 95\% interval of
$[2.6\times10^{-3},\,3.7\times10^{-2}]$ for the pooled rate.  It is retained
here as a historical work figure for that campaign, not as a rate for uniform
random contexts.}

\section{The Submask Context Family}
\label{sec:family}

\begin{definition}[The family]
A context is in the \emph{submask family} if
\[
a_4 = a_5 = v \quad\text{for some } v,\qquad e_8 = \hex{0xFFFFFFFF},\qquad
e_9 = \hex{0xFFFFFFFF}.
\]
Since $e_8 = a_4+a_8-T_2^{(8)}$ and $e_9=a_5+a_9-T_2^{(9)}$, the last two
conditions are met by choosing $a_6,a_7$ freely and then setting
$a_8 = -1-a_4+\Sig0(a_7)+\Maj(a_7,a_6,a_5)$ and
$a_9 = -1-a_5+\Sig0(a_8)+\Maj(a_8,a_7,a_6)$.  The words $v,a_6,a_7,a_{10}$
(and $a_{11}$ at $R=20$) remain free, so the family contains $2^{128}$
contexts at 19 rounds and $2^{160}$ at 20.
\end{definition}

The table solve can be written explicitly without importing notation from the
earlier paper.  In the order $j=0,1,2$, write the round-inverted words as
\begin{equation}
 W_{j+1}=a_{j+1}+B_j,\qquad W_{9+j}=A_j-a_{j+1},
 \label{eq:AB}
\end{equation}
where $B_j$ depends only on variables already solved and $A_j$ depends only on
those variables and the context.  For $j=0$, the second relation uses the
provisional $\hat W_9-a_1$; for $j=1$ it holds in the family by
Proposition~\ref{prop:triangular}; and for $j=2$ it follows directly from
$e_7=a_3+T_1^{(7)}$ in~\eqref{eq:recW}.  Substitution in~\eqref{eq:Cj} gives
\begin{equation}
 \sig0(W_{j+1})-W_{j+1}
 =W_{16+j}-\sig1(W_{14+j})-W_j-A_j-B_j=:R_j.
 \label{eq:lookupRHS}
\end{equation}
Thus every stored root $u$ of $\sig0(u)-u=R_j$ gives
$W_{j+1}=u$ and $a_{j+1}=u-B_j$.  This is the complete algebraic lookup used
at each stage; a missing index rejects the path, and a multi-root table
branches over its stored roots.

\begin{proposition}[The iteration disappears]
\label{prop:triangular}
In the submask family the right-hand side of $C_1$ does not depend on $a_3$.
Consequently, for each swept $a_0$, $a_1$, $a_2$ and $a_3$ are obtained by
three successive table lookups, $C_0\to a_1$, $C_1\to a_2$, $C_2\to a_3$, and
every triple so obtained satisfies $C_0$ (with the provisional $\hat W_9$),
$C_1$ and $C_2$ exactly.
\end{proposition}

\begin{proof}
The only route by which $a_3$ reaches $C_1$ is through $W_{10}$ in
\eqref{eq:Cj}, and by~\eqref{eq:recW}
$W_{10}=\text{const}-e_6-\Ch(e_9,e_8,e_7)$ with
$e_6 = a_2 + a_6 - \Sig0(a_5) - \Maj(a_5,a_4,a_3)$ and $e_7=a_3+T_1^{(7)}$.
The $a_3$-dependence is therefore
\[
D(a_3) = \Maj(a_5,a_4,a_3) - \Ch(e_9,e_8,\,a_3+T_1^{(7)}) .
\]
With $a_5=a_4$, $\Maj(a_4,a_4,a_3)=a_4$; with $e_9=\hex{0xFFFFFFFF}$,
$\Ch(e_9,e_8,\cdot)=e_8$.  Both terms are constants, so $D$ is constant.
$C_2$'s right-hand side depends on $a_2$ (and on context) but its own unknown
$a_3$ cancels as in the base attack.  The three constraints are therefore
triangular in $(a_1,a_2,a_3)$ and each one-root lookup either returns its
stored representative or misses; a multi-root implementation branches over
the retained representatives.
\end{proof}

Neither condition suffices alone: an isolation test on three targets gave
0--16 triangular solutions per $2^{22}$ swept $a_0$ for random contexts, 3--391
with $a_5=a_4$ only, 2--236 with $e_9=\hex{0xFFFFFFFF}$ only, and about
$1.07\times10^6$ with both, i.e.\ $0.402$ per $C_0$ survivor, which is
$0.634^2$, the probability that the two further lookups both hit.

\begin{proposition}[The consistency check collapses]
\label{prop:collapse}
In any context with $a_4=v$ and $e_8=\hex{0xFFFFFFFF}$ (the condition
$a_5=a_4$ is not needed for this part),
\begin{equation}
\varepsilon \;=\; \Maj(v,a_3,a_2) - a_3 \pmod{2^{32}} .
\label{eq:collapse}
\end{equation}
With $P = a_3\wedge\neg a_2\wedge\neg v$ and
$Q=a_2\wedge\neg a_3\wedge v$, one has $\Maj(v,a_3,a_2)=a_3-P+Q$: the first
operation clears selected one-bits and the second sets selected zero-bits, so
those two operations respectively use no borrows and no carries.  Hence
$\varepsilon=0$ if and only
if $\Maj(v,a_3,a_2)=a_3$ as bit strings, equivalently
\begin{equation}
(a_2\oplus a_3)\wedge(a_3\oplus v) = 0,
\label{eq:compact}
\end{equation}
which holds for exactly $3^{32}$ of the $2^{64}$ pairs $(a_2,a_3)$, i.e.\
with probability $(3/4)^{32}=1.004524\times10^{-4}$ for uniform independent
$a_2,a_3$, for every $v$.
\end{proposition}

\begin{proof}
With $e_8=\hex{0xFFFFFFFF}$, $\Ch(e_8,e_7,e_6)=e_7=a_3+T_1^{(7)}$ in
\eqref{eq:W9}, so the $a_1$-free part of $W_9$ is
$W_9^{\mathrm{conv}} = \text{const} + \Sig0(v) - a_5 + \Maj(v,a_3,a_2) - a_3 - T_1^{(7)}$.
Evaluating at $a_2=a_3=0$ gives the same expression with $\Maj(v,0,0)=0$ and
$a_3=0$, and subtracting yields~\eqref{eq:collapse}.  For the decomposition,
compare bits: where $v_i=0$, $\Maj$ is $a_{3,i}\wedge a_{2,i}$, which differs
from $a_{3,i}$ exactly when $a_{3,i}=1,a_{2,i}=0$, the bits of $P$; where
$v_i=1$, $\Maj$ is $a_{3,i}\vee a_{2,i}$, which differs from $a_{3,i}$ exactly
when $a_{3,i}=0,a_{2,i}=1$, the bits of $Q$.  Since $P\subseteq a_3$,
$a_3-P$ clears bits without borrows, and since $Q$ is disjoint from $a_3-P$,
adding $Q$ sets bits without carries.  Independently of this representation,
two 32-bit words have zero modular difference exactly when they are equal.
Therefore $\varepsilon=0$ iff $P=Q=0$, i.e.\ iff each bit position avoids the
two forbidden patterns; the surviving pairs number $3^{32}$.
\end{proof}

At $v=0$ condition~\eqref{eq:compact} reads $a_3\wedge\neg a_2=0$: $a_3$ is a
submask of $a_2$, which names the family.  The exactness of
Proposition~\ref{prop:collapse} was checked on 2000 random contexts and
pairs, and its probability on $2^{22}$ random pairs at $v\in\{0,
\hex{0xFFFFFFFF}, \hex{0x6A09E667}\}$, by \texttt{code/submask\_family.py}
in the repository, which needs no table.

\begin{corollary}[Yield]
\label{cor:yield}
Model the three lookups as returning words that are uniform and independent of $v$ and of one another; the count $3^{32}$ is exact for uniform independent pairs, and Section~\ref{sec:r19} tests the model.  Under it, with a single-root table of image fraction $c=0.633673$, each swept $a_0$ yields $c^3=0.2545$ triangular solutions, of which a fraction $(3/4)^{32}$ satisfy~\eqref{eq:eps}.  Every such triple is a 19-round preimage.  The
expected number of swept $a_0$ per 19-round preimage is therefore
\[
\frac{1}{c^3\,(3/4)^{32}} \;=\; 39{,}124 \;=\; 2^{15.26}.
\]
\end{corollary}

The last claim of the corollary, that the collapsed condition is
\emph{sufficient}, also follows from the construction.  Equation~\eqref{eq:recW}
makes every selected state transition hold by definition.  The constraints
$C_0,C_1,C_2$ make the expanded words $W_{16},W_{17},W_{18}$ agree with the
words fixed by the backward chain; at 20 rounds $C_3$ does the same for
$W_{19}$.  Because that chain ends at the target state, a path satisfying the
applicable constraints and the exact consistency check forward-evaluates to
the target.  As an implementation check, on 64{,}000 triangular solutions
recomputed from their message words, $C_1$ and $C_2$ held for all, $C_0$ held
exactly for those with $\varepsilon=0$, and those were preimages.

\paragraph{Why it works, in one sentence.}
The base attack's iteration converges for about one $C_0$ survivor in
$2^{16}$ rather than at the $2^{-29}$ of a random map on a 32-bit
state~\cite{Basra2026R20}, because the feedback into $C_1$ passes through a
bitwise select rather than through the whole round function; the family
removes that feedback altogether, and the same saturation that removes it
leaves the consistency residual as a difference of two bitwise functions.

\section{Implementation}
\label{sec:impl}

\paragraph{Tables.}
The published table stores one root per image value, the largest.  A random
target has a Poisson-like number of roots (the exact image fraction
$0.633673$ is slightly above $1-1/e=0.632121$, and the multiplicity
distribution is slightly under-dispersed), so keeping only one discards a
factor $1/c=1.578$ of the available continuations at each lookup.  Building
second- and third-root tables (first free slot, $u$ distinct from the roots
already stored) does not change the image fraction.  It raises the mean number
of retained roots per uniformly selected index to $0.8979$ and $0.9774$, and
the expected triangular paths per swept $a_0$ to $0.7240$ and $0.9338$;
measured values are $0.725$ (\path{logs/smoke19_2roots.log}) and $0.9338$
(\texttt{status.json} of each production run, solutions divided by swept
$a_0$).  Each table is 16~GiB; three fit an 80-GiB GPU.  Extra branches also
cost work: the measured CPU sweep rates were $2.80\times10^{5}$ and
$2.06\times10^{5}$ swept $a_0$ per second for one and two roots respectively
(\path{logs/smoke19_1root.log}, \path{logs/smoke19_2roots.log}); the GPU ratio
was not isolated.  One index in $2^{32}$ is lost to the sentinel
(the value $\hex{0x20000000}$ has the single root $\hex{0xFFFFFFFF}$, which
encodes the ``miss'' sentinel); it is documented and irrelevant to any rate.

\paragraph{Kernel.}
A first GPU implementation in eager PyTorch reached $4.47\times10^{8}$ swept
$a_0$ per second on an H100 and $2.54\times10^{8}$ on an A100, bound not by
the table gathers but by streaming 64-bit intermediates through memory.  A
fused Triton kernel that keeps the whole chain in 32-bit registers, gathers
only from the three tables, and records only the winning $a_0$ for host-side
reconstruction reaches $9.1$--$9.3\times10^{8}$ per second kernel-timed on an
A100 (\texttt{FLEET.txt}), roughly $3.6\times$ the eager kernel on the same
GPU, and sustained $8.6$--$9.0\times10^{8}$ in production including host
overhead (\texttt{status.json} of each Triton pod, swept $a_0$ divided by
elapsed time), i.e.\ $4.7$--$5.0$~s per context of $2^{32}$ swept $a_0$.  Averaged
over each pod it produces $402{,}880$--$402{,}910$ collapsed-condition
candidates per context (\texttt{status.json}, candidates divided by contexts),
against $0.93384\cdot(3/4)^{32}\cdot 2^{32}=402{,}897$ predicted from the
exact three-table yield $0.977443^3$.  Two
implementation traps are worth recording: Triton's right shift is arithmetic,
so every rotation must mask, and a 32-bit table index sign-extends into an
illegal address unless masked to 64 bits.  The kernel was validated against
the reference implementation on planted preimages before use
(Section~\ref{sec:validation}).

\paragraph{Reference implementation and verification.}
Every reported preimage is rebuilt from its recovered state into sixteen
message words and re-hashed by an implementation independent of the solver
before it is reported, and again by the standalone verifier
\texttt{code/verify\_r19.py} (standard library only, no shared code with the
attack), which prints \texttt{result: OK}.

\section{Nineteen Rounds}
\label{sec:r19}

\begin{table}[t]
\centering\footnotesize
\caption{Nineteen-round preimages from the submask family, single-root table,
CPU (numpy).  ``Solutions'' are triangular solutions; ``$\varepsilon=0$'' are
verified preimages.}
\label{tab:r19}
\begin{tabular}{@{}lrrrrr@{}}
\toprule
Run & swept $a_0$ & solutions & per survivor & $\varepsilon=0$ & $a_0$/preimage \\
\midrule
4 generated targets & $16{,}777{,}216$ & $4{,}270{,}445$ & $0.4017$ & 395 & $42{,}474$ \\
64 generated targets & $536{,}870{,}912$ & $136{,}596{,}846$ & $0.4015$ & $13{,}840$ & $38{,}791$ \\
separate implementation A$^{a}$ & $270{,}532{,}608$ & --- & --- & $6{,}915$ & $39{,}123$ \\
separate implementation B$^{a}$ & $335{,}544{,}320$ & $85{,}389{,}673$ & --- & $8{,}643$ & $38{,}823$ \\
\midrule
all-ones digest~\cite{Zaikin2026} & $4{,}194{,}304$ & $1{,}068{,}641$ & $0.4020$ & 108 & $38{,}836$ \\
all-zero digest~\cite{Davydov2025} & $4{,}194{,}304$ & $1{,}068{,}359$ & $0.4019$ & 115 & $36{,}472$ \\
\bottomrule
\end{tabular}
\begin{minipage}{0.95\textwidth}\footnotesize\vspace{2pt}
$^{a}$ Twenty-four contexts, twenty-four generated targets and twelve values of $v$, run
from separately written code in an automated review pass that used the same
specification and repository evidence (see the disclosure); 6{,}915 of 6{,}915
collapsed-condition candidates verified as preimages; the second replication,
$8{,}643$ of $8{,}643$.  Derived value: $39{,}124$.  The four-target row is the
run reproduced bit for bit in \texttt{reviewers/reports.json}; the other rows
are in \path{logs/r19_big.log} and \path{logs/r19_fixed_targets.log}.
\end{minipage}
\end{table}

Table~\ref{tab:r19} summarises the 19-round measurements.  The 64-target run took 151.7~s of wall time in one unpinned NumPy process
($3.5\times10^{6}$ swept $a_0$ per wall-second); pinned to a single core the
rate is $3.4\times10^{6}$ swept $a_0$ per core-second
(\texttt{reviewers/reports.json}, \texttt{taskset} with
\texttt{OMP\_NUM\_THREADS=1}).  Combining this throughput with the measured
yield gives a marginal estimate of about 12~ms of one core per preimage after
the reusable table has been built; it is not an observed cold-start latency.
The measured collapse rate,
$1.0132\times10^{-4}$ per solution on the 64-target run and
$1.0122\times10^{-4}\pm1.1\times10^{-6}$ on a second, larger reviewer
replication ($3.36\times10^{8}$ swept $a_0$), agrees with $(3/4)^{32}=
1.0045\times10^{-4}$; an absolute yield audit on a separate corpus ($4.70\times10^{9}$ swept $a_0$
over 100 contexts at 19 and 20 rounds, $119{,}948$ collapsed-condition
candidates, every 19-round one a verified preimage), observed divided by the
prediction of Corollary~\ref{cor:yield} with the measured table fill
$0.63375$, gives $0.9986\pm0.0029$, so any
hidden gain or loss lies within a hundredth of a bit.  No material target
effect is resolved in these samples: the all-ones and all-zero benchmark
digests have observed yields near those of the generated targets.

\paragraph{Comparison.}
Table~\ref{tab:compare} keeps unlike hardware and cost measures explicit.  Two cautions apply.  The
new cost is a \emph{marginal} cost: both our attacks presuppose the same
one-time 16-GiB table, which takes about 215 recorded CPU core-seconds to build, so a
cold-start cost for a single preimage is dominated by the table.  And the
baseline of~\cite{Basra2026R19} is a mixed screened/control campaign rather
than a hardware-matched benchmark.  Its pooled work figure, $2^{38.3}$ swept
$a_0$, divided by the present measured $2^{15.25}$, is a work-count ratio of
about $2^{23}$ for those two implementations.  We do not turn the CPU and GPU
timings in Table~\ref{tab:compare} into a controlled speed ratio.  Zaikin's
18~h~33~min is the duration of one successful run; the stated one-in-eight
acceptance probability of its intermediate solutions does not by itself
establish an eightfold mean runtime.

\begin{table}[t]
\centering\footnotesize
\caption{Nineteen-round preimages of the SHA-256 compression function,
standard IV, all 256 bits fixed.  Costs are per preimage.}
\label{tab:compare}
\begin{tabular}{@{}l>{\raggedright\arraybackslash}p{3.3cm}>{\raggedright\arraybackslash}p{3.4cm}>{\raggedright\arraybackslash}p{3.7cm}@{}}
\toprule
 & Zaikin~\cite{Zaikin2026} & \cite{Basra2026R19} & this work \\
\midrule
Targets & one fixed (all-ones) & six generated; one fixed & generated and fixed benchmarks \\
Method & tuned Kissat inside Cube-and-Conquer & random context, fixed-point iteration & submask family, three lookups \\
Cost & 18\,h\,33\,min on 192 cores$^{a}$ & $2^{38.3}$ swept $a_0$ (pooled mixed-arm campaign); about 2 GPU-hours & $2^{15.26}$ swept $a_0$; about 12 ms marginal CPU time \\
Preimages verified & 1 & 7 & more than 20{,}000$^{b}$ \\
\bottomrule
\end{tabular}
\begin{minipage}{0.95\textwidth}\footnotesize\vspace{2pt}
$^{a}$ About 3{,}560 core-hours for the reported successful run, obtained from
an intermediate problem whose solutions are 19-round preimages with
probability $1/8$.  This single duration is not an estimate of a mean.  Direct
attempts on the 19-round instance found nothing in seven days on 192 cores
(Zaikin's account).
$^{b}$ Counts are in the logs; verified blocks with their targets are printed
in \path{logs/r19_scale.log} (33), \path{logs/r19_fixed_targets.log} (4),
\path{logs/r19_variants.log} and \path{logs/r19_degenerate.log} (2 each) and
the \path{logs/smoke19_*.log} files.
\end{minipage}
\end{table}

\section{Twenty Rounds}
\label{sec:r20}

At $R=20$ the backward chain fixes $a_{12},\dots,a_{19}$, the context is
$a_4,\dots,a_{11}$, and a fourth constraint $C_3$ on $W_{19}$ appears.  The
submask family transplants unchanged: Propositions~\ref{prop:triangular}
and~\ref{prop:collapse} involve only $a_4,\dots,a_9$ and $W_9$, and the
measured rates at 20 rounds are $0.4015$ solutions per $C_0$ survivor on 40
targets of $2^{26}$ swept $a_0$ each (\texttt{logs/c3\_uniformity.log}) and a
collapse rate of $1.0049\times10^{-4}$ per solution on $8.4\times10^{10}$
swept $a_0$ (\texttt{logs/c3\_scaling.log}).  Everything therefore reduces to the residual of $C_3$,
\begin{equation}
c_3 \;=\; \sig0(W_4) + W_3 - \kappa_3 ,\qquad
\kappa_3 = W_{19}-\sig1(W_{17})-W_{12},
\label{eq:c3}
\end{equation}
which must vanish, where $W_3,W_4$ depend on $a_0,\dots,a_4$ and $\kappa_3$ on the
context and digest only.

\subsection{Measured behaviour of the fourth residual}
\label{sec:c3}

We tested the distribution of $c_3$ on candidates that satisfy everything
else.
\begin{itemize}
\item A chi-squared test on 68{,}414 candidates, on the low and the high
  4, 6, 8, 10 and 12 bits: every $|z|\le 2.1$ (largest: high 10 bits, $z=+2.10$).
\item A scaling ladder on $8.4\times10^{10}$ swept $a_0$ ($2.14\times10^{6}$
  candidates): the number with the low $k$ bits of $c_3$ zero, against the
  uniform prediction, has ratio $1.01$ at $k=8$, $0.96$ at 12, $1.04$ at 16,
  $0.86$ at 18 (7 against 8.2).
\item The GPU campaign's own counters
  (\texttt{gpu\_run\_snapshot/h100\_status.json}, $2.96\times10^{9}$
  candidates): the counts $n_k$ are nested.  Under the uniform-residual null,
  $n_{k+1}$ is conditionally Binomial$(n_k,\tfrac12)$, but the observed ratios
  do not thereby become unconditionally independent because their denominators
  are random.  A score formed from the conditional transitions over bits
  20--28 gives $2{,}654$ successes of $5{,}370$ opportunities on the H100
  (expected $2{,}685$, $z=-0.85$).  The most extreme individual transition,
  selected after examining the eight steps, has exact unadjusted one-sided
  $p=0.00843$ and Bonferroni-adjusted $p=0.0674$.  The two A100 runs give
  $z=+1.47$ and $z=+1.71$ over the same transitions.  These are descriptive
  large-sample score calculations, not independent-binomial replications.
\item A dedicated test on the all-ones digest itself, against generated digests
  as a matched control (\texttt{logs/ones\_obstruction.log};
  $3.95\times10^{6}$ candidates in each arm): no deficit that deepens with $k$ appears in either arm: every cell of the
  all-ones arm from 8 bits on is within $1.3\sigma$ of expectation below 18
  bits, and the only cell outside $2.2\sigma$ on the low side anywhere is the
  control arm's 6-bit count ($61{,}030$ against $61{,}786$, $z=-3.0$, a $1.2\%$
  shortfall absent at 8 bits and beyond).  The all-ones arm runs high between
  18 and 20 bits (at 20 bits, 11 observed against 3.8 expected, $z=+3.7$;
  conditional halving over steps 16--21 $z=+3.3$) and returns to expectation
  at 22 (1 against 0.9); the statistic was examined after each of six rounds
  with a widening window, so we read the surplus as noise and the test as
  excluding only an obstruction, which is what it was built for.
\end{itemize}
The deepest cell with useful occupancy is the 24-bit cell (174 candidates
against 176.5 expected on the H100 run).  Counts at 28 bits are too sparse to
bound an atom at exact zero without distributional assumptions.  Thus
$P(c_3=0)=2^{-32}$ is a working extrapolation, not a measured exact-zero rate.
The two observed hits in Section~\ref{sec:hits} are consistent with it but do
not estimate it precisely.

\paragraph{Why we find no steering.}
For a fixed candidate $(a_0,\dots,a_3)$ the digest and the context words
$a_5,\dots,a_{11}$ enter~\eqref{eq:c3} only through the additive constant
$\kappa_3$; the word $v=a_4$ enters $W_4$ additively before $\sig0$, and the
context constants also determine which $(a_1,a_2,a_3)$ the three lookups
return.
Writing $S=\sig0(W_4)+W_3$, a sum modulo $2^{32}$, the constant contributes
only a phase, $\varphi_{c_3}(\beta)=\varphi_{S}(\beta)\,e^{-2\pi i\beta
\kappa_3/2^{32}}$, and if $W_3$ and $\sig0(W_4)$ were independent $\varphi_S$
would factorise as $\varphi_{W_3}\varphi_{\sig0(W_4)}$; the two share
$a_1,a_2,a_3$, so we treat the factorisation as a model and test it on 40
contexts of about $1{,}700$ candidates each.  The largest-root policy gives
$W_3$ a visible first harmonic, $|\varphi_{W_3}(1)|=0.155$, while the observed
$|\varphi_{\sig0(W_4)}(1)|=0.025$ and $|\varphi_{c_3}(1)|=0.026$ are near the
per-context sampling scale $1/\sqrt n=0.024$.  An $8.4\times10^6$-frequency
scan found no statistic above its uniform-model reference maximum; the same
scan detects the known $W_3$ bias.  Separate mutual-information tests against
$a_0$, $v$, context identifiers and bit fields, and random linear
approximations from $a_0$, likewise produced maxima comparable with their
shuffle or noise controls.

These tests concern the sampled context distribution, the implemented root
policies and the statistics just named.  They found no usable steering signal
at their detection limits.  Marginal and pairwise tests do not establish the
joint factorisation above, exclude sparse exact-zero effects, or rule out an
unexamined table policy or context family.

\subsection{Cost}
\label{sec:cost20}

Under the working assumptions that the retained lookup branches follow their
measured mean, the collapsed pairs have the uniform-pair rate of
Proposition~\ref{prop:collapse}, and $c_3=0$ has probability $2^{-32}$, the
probability of a 20-round preimage per swept $a_0$ is
$0.9337\cdot(3/4)^{32}\cdot 2^{-32}=2.18\times10^{-14}$ with three-root tables,
i.e.\ $2^{45.4}$ swept $a_0$ per preimage ($2^{47.3}$ with one root).  At
$8.5\times10^{8}$ per second on an A100 this model gives a mean of about
15~h.  If successive candidates have a constant, independent success
probability, the waiting-time model is approximately exponential; that is an
additional modelling assumption, not a consequence of the two observed
successes.  We do not extrapolate an H100 time because the fused production
kernel was not run on that device.

\subsection{The computed preimages}
\label{sec:hits}

\begin{table}[t]
\centering\small
\caption{The recorded 20-round production runs.  Every run built three root
tables on its device.  ``Model mean'' is the number of collapsed candidates
times the assumed $2^{-32}$ fourth-constraint rate.}
\label{tab:campaign}
\begin{tabular}{@{}lllrrrl@{}}
\toprule
GPU & kernel & target & contexts & swept $a_0$ & model mean & outcome \\
\midrule
H100 & PyTorch & generated digest & $7{,}351$ & $3.16\times10^{13}$ & 0.690 & none \\
A100 & PyTorch & generated digest & $1{,}839$ & $7.90\times10^{12}$ & 0.173 & none \\
A100 & PyTorch & generated digest & $1{,}836$ & $7.89\times10^{12}$ & 0.172 & none \\
A100 & Triton & generated digest & $1{,}765$ & $7.58\times10^{12}$ & 0.166 & \textbf{preimage}, 2.45 h \\
A100 & Triton & all-ones & $12{,}203$ & $5.24\times10^{13}$ & 1.145 & \textbf{preimage}, 16.16 h \\
A100 & Triton & all-ones & $12{,}075$ & $5.19\times10^{13}$ & 1.133 & none (stopped) \\
\midrule
total & & & & $1.59\times10^{14}$ & 3.48 & 2 preimages \\
\bottomrule
\end{tabular}
\end{table}

Table~\ref{tab:campaign} lists every main production run (omitting about 48 contexts,
$2\times10^{11}$ swept $a_0$, on generated targets from the two all-ones runs
before they were retargeted; \texttt{FLEET.txt}).  The six status records sum
to 71.7 GPU-hours of solving time.  They produced two preimages against a
model mean of 3.48.  Because successful jobs stopped on a hit and related jobs
were then stopped manually, this campaign was not a fixed-exposure experiment;
we therefore report the comparison descriptively and do not attach a Poisson
significance test to it.

\paragraph{Preimage of a solver-generated digest.}
The target is the 20-round digest of a random 55-byte message drawn by the
solver, so a preimage was known to exist; the block found is a different one.
\begin{center}\ttfamily\small
target: d962ca30 635f9b74 ac6c8c12 43a1a9cf 800e81bc 05f1d2e4 0764c68c 795f7388\\[3pt]
block:\phantom{ } a36f4238 f2c9204e b5b6653b 070401f7 5928e4f3 fe766be2 52026907 f7ee1812\\
\phantom{block: }01344603 ea505012 86cbe6cb 6ce73d0b 5c91de6a 43355e3b ff3d5e88 c1ad0b54
\end{center}
Found 2026-09-04 in context 1{,}765 after 2.45~h on one A100
($7.58\times10^{12}$ swept $a_0$).  The recovered state has
$a_4=a_5=\hex{0x799f3424}$, $e_8=e_9=\hex{0xFFFFFFFF}$ and satisfies~\eqref{eq:compact}.

\paragraph{Preimage of the all-ones digest.}
Target $\hex{ff}\cdots\hex{ff}$ (256 one-bits), the digest Zaikin~\cite{Zaikin2026}
inverted at 19 rounds and chose because, for an irregular output, ``it may be hard to justify
this choice and prove that a random pair of input and output was not picked''
(\cite{Zaikin2026}, \S5.1); no 20-round preimage of it
was known to exist:
\begin{center}\ttfamily\small
0b0187bf b9aea692 b66effa5 087dca3a 0caea827 1d2f9916 739a224e a87c0eae\\
7f9ef4a7 b318a7de a8848c61 a7a141a4 11ac114b 952299be 97aaa67f c6acfe57
\end{center}
Found 2026-09-05 in context 12{,}203 after 16.16~h on one A100
($5.24\times10^{13}$ swept $a_0$); a second A100 on the same digest swept
$5.19\times10^{13}$ in 12{,}075
contexts without a hit before being stopped.  The recovered state has $a_4=a_5=\hex{0x19eb97f7}$,
$e_8=e_9=\hex{0xFFFFFFFF}$ and satisfies~\eqref{eq:compact}.

Both blocks verify with \texttt{code/verify\_r19.py --rounds 20}, fail with
\texttt{--rounds 19} and \texttt{--rounds 21}, and reproduce the digest under
a separately written from-scratch SHA-256.  Appendix~\ref{app:witness} gives the
commands.

\paragraph{The same digest, one more round.}
Zaikin inverted the all-ones digest at 19 rounds in 18~h~33~min on 192 CPU
cores.  We use the same fixed target at 20 rounds, but the unlike algorithms,
hardware and single successful observations are not a controlled performance
comparison.

\subsection{Validation of the pipeline}
\label{sec:validation}

Because no computed 20-round preimage was available to us before the
campaign, the fourth-constraint code path had never been exercised on a
genuine input, and an error in it
would leave the residual looking uniform while silently discarding real hits.
We therefore \emph{planted} 20-round preimages: blocks constructed so that
their internal state lies in the family and satisfies~\eqref{eq:compact}, then
hashed for 20 rounds.  Run over a window containing the true $a_0$, the
reference implementation found every plant whose table roots were stored (10
of 16 with two tables; a separate automated replication found 54 of 200 with one
table and 145 of 200 with two, against the expected recovery fractions
$c^3=0.25$ and $0.8979^3=0.72$) and never a wrong one.  The Triton kernel was
then checked against the reference on 12 plants with two and three tables,
agreeing on every hit, and an interpreter-mode execution compared all 27
root-combination paths against the reference on identical inputs.

\section{Twenty-One-Round Limitations}
\label{sec:barrier}

Each round beyond 19 adds one schedule constraint.  The companion
note~\cite{Basra2026R20} found that, for its tested random-context
parameterisations, absorbing a fifth unknown creates a dependency cycle whose
largest measured edge is $a_4\to C_0$ through $\Sig0(a_4)$ in $W_9$.  The
submask family does not cancel that edge; it removes a different dependency
($a_3\to C_1$) and tests $C_3$ directly.  We report the finite searches below
as limitations of the present construction, not as an impossibility or
optimality result.

\paragraph{Context families tested.}
The saturation mechanism of Proposition~\ref{prop:triangular} clears a
dependency when it passes through $\Maj$ or $\Ch$ with a context word among
the arguments.  It supplies no direct lever for a dependency through $\Sig0$
or $\sig0$ of the unknown itself.  For $C_3$ to absorb $a_4$ in this triangular
arrangement one would need $C_1$ free of
$a_3$ and of $a_4$, and $C_2$ free of $a_4$.  Table~\ref{tab:frames} reports,
for six condition families, which of these hold, measured by perturbing each
unknown and observing the constraint targets.

\begin{table}[t]
\centering\small
\caption{Observed dependency removals for six tested context families at 20
rounds (40 random instances each, measured by
\texttt{scripts/frame\_search2.py}; output in \texttt{logs/frame\_search.log}).}
\label{tab:frames}
\begin{tabular}{@{}lccc@{}}
\toprule
Context conditions & $C_1$ free of $a_3$ & $C_1$ free of $a_4$ & $C_2$ free of $a_4$ \\
\midrule
random & no & no & no \\
$a_4=a_5$, $e_8=e_9=\hex{-1}$ (this family) & \textbf{yes} & no & no \\
$\ldots$ and $a_6=a_5$ & \textbf{yes} & no & no \\
$a_5=a_6$, $e_9=e_{10}=\hex{-1}$ (shifted) & no & no & \textbf{yes} \\
$a_4=a_5=a_6$, $e_8=e_9=e_{10}=\hex{-1}$ & \textbf{yes} & no & \textbf{yes} \\
$a_4=\dots=a_7$, $e_8=\dots=e_{11}=\hex{-1}$ & \textbf{yes} & no & \textbf{yes} \\
\bottomrule
\end{tabular}
\end{table}

The middle column remains ``no'' in these six families.  With $a_5=a_4$ the $a_4$-dependence of $C_1$
is $\Sig0(a_4)$ inside $T_1^{(6)}$; without it, $a_3$ leaks back in through
$\Maj(a_5,a_4,a_3)$.  A symbolic search over 1{,}485 context-condition families
from the implemented grammar (\path{analysis/scripts/symdep/}) found none
that clears both routes.  The collapse of Proposition~\ref{prop:collapse} survives the
tighter families (collapse rates $1.016\times10^{-4}$ and
$0.987\times10^{-4}$, \texttt{logs/frame\_search.log}), but they gave no gain
under the tested solve order.

\paragraph{A recurring heavy edge.}
The broader finite search localises the failure within its representation.
Two recurring facts were observed at 20 rounds
(\path{analysis/elimination_search/}).  First, an
\emph{offset law}: every state word $a_k$ enters the recovered message word
$W_{k+5}$ through $\Sig0(a_k)$, carried by the term
$e_{k+1}=a_{k-3}+a_{k+1}-\Sig0(a_k)-\Maj(a_k,a_{k-1},a_{k-2})$ subtracted
inside~\eqref{eq:recW}, and through $\Sig1(e_{k+4})$ with $e_{k+4}=a_k+\text{const}$;
both are one-input functions of $a_k$ itself, so the tested saturation rule
does not apply to them.  Second, $a_1$, $a_2$ and $a_3$ had fixed roles in
every one of
$11{,}086$ legal context configurations examined, $a_1$ is heavy in $C_1$, $C_2$
and $C_3$ and can be absorbed only by $C_0$, and likewise $a_2$ only by $C_1$
and $a_3$ only by $C_2$.  Within this arrangement, a fourth absorbed word
would therefore need to be light in
$C_0$, $C_1$ and $C_2$ and absorbable by $C_3$, and the offset law puts
$\Sig0(a_k)$ heavily into at least one of $C_0$, $C_1$, $C_2$ for every
$k=4,\dots,11$.

The closest miss is $a_5$, and it localises the barrier to one wire.  With
$a_6=\neg a_4$, $a_7=a_6$ and $e_8=\hex{0xFFFFFFFF}$ the word $a_5$ is
\emph{exactly} absent from $C_0$ and from $C_2$, and with $e_{11}=0$ it enters
$C_3$ exactly linearly; every path but one is cut.  The one that remains is
$a_5\to C_1$ through $-\Sig0(a_5)$ inside $e_6=a_2+a_6-\Sig0(a_5)-\Maj(a_5,a_4,a_3)$
in $W_{10}$.  With $a_0,\dots,a_3$ fixed, the $C_1$ lookup target depends on
$a_5$ through exactly four atoms,
$-\Sig0(a_5)-\Maj(a_5,a_4,a_3)+\Sig1(a_5+c_9)+\Ch(a_5+c_9,e_8,e_7)$, where
$c_9=a_9-T_2^{(9)}$ is the only context lever besides flattening the
$\Maj$ and $\Ch$ terms (worth one to two bits).  The kernel
$\Sig1(x+c_9)-\Sig0(x)$ changes at least $10.3$ bits per flipped bit of $x$
for every $c_9$ in a scan of $201{,}058$ values including every sparse one.
A falsification test over about $1.6\times10^{6}$ legal context conditions
outside the standard menu (arbitrary constants for $e_8,\dots,e_{11}$ and the
context words; rotational, XOR, additive, $\Sigma$ and GF(2)-linear ties;
a symbolic search for a cancelling second occurrence; a full $2^{32}$ scan of
the frozen-$a_5$ residual; and a reduced-width model exhaustive over all
$2^{25}$ contexts at five bits and checked at 6, 8, 12 and 32) found a global
minimum of $10.86$ bits per flipped bit, against the 2-bit threshold at which
absorption becomes possible and a control of $12.3$--$13.1$ bits on the
$a_4\to C_0$ edge; the frozen residual had random-map-like measured
collision statistics (empirical zero rate $2\cdot2^{-32}$ and no detected
affine output bit)
(\path{analysis/elimination_search/falsify_a5/}).  The obstruction to a fourth
absorption in these experiments is thus localised to this single edge.  It is
heavy for the same reason as the $a_4\to C_0$ edge of
\cite{Basra2026R20}: $\Sig0$ of the unknown itself.  This is a negative result
for the stated representation, condition grammar and solve order, not for all
table-based or algebraic methods.

\paragraph{Joint-residual measurement.}
On 214{,}583 freshly generated 21-round candidates (84 targets,
$8.4\times10^{9}$ swept $a_0$; \path{analysis/scripts/symdep/gen.log})
satisfying everything up to the collapsed condition, the residuals of $C_3$ (at $W_{19}$) and $C_4$ (at
$W_{20}$) were individually consistent with the tested uniformity statistics,
and the estimated pairwise mutual information was below $10^{-5}$ bits.  This
does not prove joint independence or bound sparse dependence at exact zero.
If both residuals are modelled as independent uniform 32-bit words, 21 rounds
adds a factor $2^{64}$ relative to the 19-round candidate generator (or
$2^{32}$ relative to the 20-round model), for about $2^{77}$ swept $a_0$ with
three-root tables.  This is a conditional estimate, not a lower bound.

\paragraph{Exploratory corpus checks.}
We mined 51{,}268 verified 19-round preimages and 68{,}680 20-round
candidates for structure not forced by the construction: per-bit biases,
GF(2)-linear and affine relations, context effects, neighbour structure in
$a_0$, residue-class structure in $c_3$, and the entropy accounting of
Proposition~\ref{prop:collapse}.  The 32 per-position marginal entropies of $(a_2,a_3)$ sum to $50.7189$ bits
against the predicted $32\log_2 3=50.7188$ (each position
$1.58496\pm0.00003$), and the 496 estimated pairwise mutual informations between
positions sum to $0.001$ bits.  A sum of marginal entropies is not the joint
entropy, and pairwise tests do not exclude higher-order dependence; these
numbers show only agreement at first and second order.  The GF(2) rank deficit is exactly the three carry-free least-significant
bit relations the construction implies; and no other test departed from its
reference distribution.  Seven separate automated analyses (see the
disclosure) proposed eight candidate levers, and a separate refutation pass
rejected every one within the tested representations.  One methodological note: a corpus that
draws $a_0$ with replacement contains duplicate rows that produce a spurious
$69\sigma$ ``collision'' signal in $c_3$ until deduplicated.

\section{Prior Work and Priority}
\label{sec:priority}

\paragraph{Computed preimages of round-reduced SHA-256.}
16 rounds by SAT in 2012~\cite{Homsirikamol2012}; 17 and 18 rounds for the
all-zero digest, 19 only in a weakened form (at most 31 of the 256 digest bits fixed, to
zero: 32 instances at one-bit steps, all but the 32-bit one solved), and
19-round collisions, by Davydov et~al.~\cite{Davydov2025}; 19 rounds for the
all-ones digest by Zaikin~\cite{Zaikin2026}, extending~\cite{Zaikin2024CP,
Zaikin2024JAIR}; and seven target-only 19-round examples from our earlier
construction, six for image-derived targets and one for the same
all-ones target~\cite{Basra2026R19}.  The repository records four automated
search passes, completed in September 2026, over the IACR ePrint archive and
proceedings, dblp, OpenAlex, arXiv, selected SAT, CP and AI venues, relevant
repositories, and some Russian- and Chinese-language sources.  Those searches
found no prior published, computed witness at 20 or more rounds in the same
model: standard IV, feed-forward, one unrestricted block and all 256 output
bits fixed.  This is a dated literature finding, not a proof of priority.  The
search did not systematically enumerate CP~2026, IJCAI~2026, AAAI~2026,
doctoral theses, or SAT-competition logs.

\paragraph{Theoretical preimage attacks.}
The meet-in-the-middle line~\cite{IsobeShibutani2009,AokiGuo2009,KRS2012}
reaches 45 steps at $2^{255.5}$; Guo et~al.~\cite{GuoLiLiuWangZhang2026}
extend it to 46 steps at $2^{254.5}$ and 47 at $2^{255.6}$ (pseudo-preimages
at $2^{253.2}$), improve 43--45 steps, take SHA-512 to 51 steps, and describe
the preceding decade as stagnant.  None of these attacks has been executed;
their preimages are obtained from pseudo-preimages by the generic conversion
and pay for padding.  Rounds and steps coincide for SHA-256 in both
literatures, so no prior result becomes twenty under a different convention.
The citation runs one way: the SAT line quotes the meet-in-the-middle
bounds~\cite{Zaikin2026,Davydov2025}, but \cite{GuoLiLiuWangZhang2026}
contains no SAT-based preimage work.

\paragraph{Computed partial results that must not be conflated.}
The artefact accompanying~\cite{GuoLiLiuWangZhang2026}
(the repository \path{github.com/desert-oasis/dual-sync-MITM},
directory \path{SHA256-44_partial_pseudo_preimage}) computes messages for 44-step
SHA-256 whose output has 32 or 40 of 256 bits fixed to zero, from a free
chaining value: weakened on three axes at once.  Bouam et~al.~\cite{Bouam2021} exhibit three messages whose \emph{full} SHA-256 digests XOR to zero on
their first 128 bits (a 3XOR), a brute-force computation of at least
$8.7\times10^{22}$ integer operations by the authors' count; that is a collision-type result on half the output, not a preimage
of any bits.  The GraphBench benchmark~\cite{Stoll2025}, building on the
encodings of~\cite{Nejati2017,Alamgir2024}, emits instances labelled as
SHA-256 preimages at 16--30 rounds, but with up to 384 of the 512 input bits
fixed, and no round count has been claimed from them.

\paragraph{The levers.}
Choice-function absorption is the message-stealing tool of
\cite{AokiGuo2009}; majority degeneracy is formalised in
\cite{GuoLiLiuWangZhang2026}; our companion note~\cite{Basra2026R20}
records both on a different constraint.  The pairing of two context words,
the application to the iteration-carrying constraint, and the bitwise
collapse of the modular consistency residual are, to our knowledge, new.

\section{Relation to Our Earlier Papers}
\label{sec:relation}

This paper supersedes the cost figures of~\cite{Basra2026R19} and the
headline of~\cite{Basra2026R20}, and we say so precisely rather than
silently.

\begin{itemize}
\item \cite{Basra2026R19} reports about one preimage per 80 attempts of $2^{32}$
  swept $a_0$, i.e.\ about $2^{38.3}$ swept $a_0$, when the screened and
  control arms of its campaign are pooled.  All three events were in the
  screened arm; the control arm had none.  The figure describes that mixed
  campaign, not uniformly sampled contexts.  The submask family is a distinct
  construction with the same table and backward chain, measured at
  $2^{15.25}$.
\item \cite{Basra2026R20} estimates a factor $2^{32}$ for the fourth
  constraint under a uniform-residual model and locates the tested
  dependency in $\Sig0(a_4)$ inside $W_9$.  The exact residual equation and
  the finite negative searches remain useful; their extrapolation is a model,
  not a bound (Section~\ref{sec:barrier}).  What changes
  is the per-context cost of the first three constraints, which the family
  reduces from a lossy iteration to three lookups.  The note's own caveat,
  that its negative results are properties of the parameterisations tried,
  is exactly what applied: its section on a context family that removes
  $a_2$ from $C_0$ applies the same two levers to $C_0$ and records the result
  as a structural observation rather than an attack step, because the $\Maj$
  condition it chose tied a context word to an unknown; applying them to $C_1$
  with a pairing on two context words is the step it did not take.
\item The screening campaign of~\cite{Basra2026R19} reports a $3.41\times$
  lift in a downstream counter.  That figure is trajectory-counted;
  duplicates are $0.03\%$ of converged trajectories, and identity-level deduplication can move it only
  within $[2.42\times,3.44\times]$.  It is now of historical interest only.
\end{itemize}

\section{Conclusion}
\label{sec:conclusion}

Imposing three conditions on the context words removes the iteration from the
19-round construction and reduces its final consistency check to a bitwise
condition.  The measured work is $2^{15.25}$ swept values per preimage, within
one per cent of the yield model; the corresponding 12-ms CPU figure is a
marginal estimate after reusable preprocessing.  At 20 rounds the exact
fourth constraint remains.  Modelling its residual as uniform gives
$2^{45.4}$ swept values and about fifteen A100 hours per preimage; two
computed examples are exhibited, including one for the all-ones target.
The six production records total 71.7 GPU-hours, while the two successful
runs took 2.45 and 16.16 hours.

To our knowledge, as of September 2026 and after the search described in
Section~\ref{sec:priority}, no computed preimage of the SHA-256 compression
function beyond 19 rounds had previously been published in the same
full-output, fixed-IV model.  At 21 rounds the tested residuals and context
families support, but do not prove, an independent-uniform model for two
32-bit constraints.  No impossibility or lower-bound claim follows.  Full and
padded SHA-256 are unaffected.

\section*{Acknowledgements and disclosure}
The derivations, code, experiments, verification and text of this paper were
developed with substantial assistance from Anthropic's Claude, used through
Claude Code under the author's direction, including experiment design, code,
analysis and drafting.  OpenAI Codex assisted with the present evidence audit,
claim qualification and manuscript revision.  The supplied blocks were
forward-checked by the standalone verifier and by several separately written
SHA-256 implementations.  Other automated review sessions variously
reimplemented parts of the construction or inspected shared logs and code;
their reports are evidence of those specific checks, not human peer review or
blanket independent replication.  The artefacts are listed in
Appendix~\ref{app:repro}.

\begin{thebibliography}{10}

\bibitem{Basra2026R19}
K.~S.~Basra.
\newblock A practical preimage attack on the 19-round {SHA-256} compression
  function via a global $\sigma_0$-difference table.
\newblock Manuscript, 2026.  Source, code and witnesses at
  \url{https://github.com/kamb-code/sha256-r19-preimage}.

\bibitem{Basra2026R20}
K.~S.~Basra.
\newblock The 20-round barrier for table-based preimage attacks on
  \mbox{{SHA-256}}.
\newblock Companion note, 2026.

\bibitem{FIPS180}
National Institute of Standards and Technology.
\newblock \emph{{FIPS PUB 180-4}: Secure Hash Standard ({SHS})}, 2015.

\bibitem{IsobeShibutani2009}
T.~Isobe and K.~Shibutani.
\newblock Preimage attacks on reduced {Tiger} and {SHA-2}.
\newblock In \emph{Fast Software Encryption (FSE 2009)}, LNCS 5665,
  pp.~139--155. Springer, 2009.

\bibitem{AokiGuo2009}
K.~Aoki, J.~Guo, K.~Matusiewicz, Y.~Sasaki, and L.~Wang.
\newblock Preimages for step-reduced {SHA-2}.
\newblock In \emph{Advances in Cryptology -- ASIACRYPT 2009}, LNCS 5912,
  pp.~578--597. Springer, 2009.

\bibitem{KRS2012}
D.~Khovratovich, C.~Rechberger, and A.~Savelieva.
\newblock Bicliques for preimages: attacks on {Skein-512} and the {SHA-2}
  family.
\newblock In \emph{Fast Software Encryption (FSE 2012)}, LNCS 7549,
  pp.~244--263. Springer, 2012.

\bibitem{GuoLiLiuWangZhang2026}
J.~Guo, H.~Li, M.~Liu, S.~Wang, and T.~Zhang.
\newblock Dual-syncopation meet-in-the-middle attacks: new results on
  {SHA-2} and {MD5}.
\newblock In \emph{Advances in Cryptology -- EUROCRYPT 2026}, LNCS 16546, pp.~121--151. Springer, 2026.
  DOI 10.1007/978-3-032-25333-0\_5.  Cryptology ePrint Archive, Report
  2026/353.

\bibitem{Homsirikamol2012}
E.~Homsirikamol, P.~Morawiecki, M.~Rogawski, and M.~Srebrny.
\newblock Security margin evaluation of {SHA-3} contest finalists through
  {SAT}-based attacks.
\newblock In \emph{Computer Information Systems and Industrial Management
  (CISIM 2012)}, LNCS 7564, pp.~56--67. Springer, 2012.
  DOI 10.1007/978-3-642-33260-9\_4.

\bibitem{Davydov2025}
V.~V.~Davydov, M.~D.~Pikhtovnikov, A.~P.~Kiryanova, and O.~S.~Zaikin.
\newblock Analysis of the cryptographic strength of the {SHA-256} hash
  function using the {SAT} approach.
\newblock \emph{Scientific and Technical Journal of Information
  Technologies, Mechanics and Optics}, 25(3):428--437, 2025.
  DOI 10.17586/2226-1494-2025-25-3-428-437.

\bibitem{Zaikin2024CP}
O.~Zaikin.
\newblock Inverting step-reduced {SHA-1} and {MD5} by parameterized {SAT}
  solvers.
\newblock In \emph{30th International Conference on Principles and Practice
  of Constraint Programming (CP 2024)}, LIPIcs 307, pp.~31:1--31:19, 2024.
  DOI 10.4230/LIPIcs.CP.2024.31.

\bibitem{Zaikin2024JAIR}
O.~Zaikin.
\newblock Inverting cryptographic hash functions via {Cube-and-Conquer}.
\newblock \emph{Journal of Artificial Intelligence Research}, 81:359--399,
  2024.  DOI 10.1613/jair.1.15244.

\bibitem{Zaikin2026}
O.~Zaikin.
\newblock Preimage attacks on round-reduced {MD5}, {SHA-1}, and {SHA-256}
  using parameterized {SAT} solver.
\newblock \emph{Constraints}, 2026 (received 28 March 2025, accepted 16 December
  2025, published online 23 January 2026).  DOI 10.1007/s10601-025-09383-0.

\bibitem{Bouam2021}
M.~Bouam, C.~Bouillaguet, C.~Delaplace, and C.~No\^us.
\newblock Computational records with aging hardware: controlling half the
  output of {SHA-256}.
\newblock \emph{Parallel Computing}, 106:102804, 2021.
  DOI 10.1016/j.parco.2021.102804.  Cryptology ePrint Archive, Report
  2021/886.

\bibitem{Nejati2017}
S.~Nejati, J.~H.~Liang, C.~Gebotys, K.~Czarnecki, and V.~Ganesh.
\newblock Adaptive restart and {CEGAR}-based solver for inverting
  cryptographic hash functions.
\newblock In \emph{Verified Software: Theories, Tools, and Experiments
  (VSTTE 2017)}, LNCS 10712, pp.~120--131. Springer, 2017.
  DOI 10.1007/978-3-319-72308-2\_8.

\bibitem{Alamgir2024}
N.~Alamgir, S.~Nejati, and C.~Bright.
\newblock {SHA-256} collision attack with programmatic {SAT}.
\newblock In \emph{Joint Proceedings of the 9th Workshop on Practical
  Aspects of Automated Reasoning (PAAR) and the 9th Satisfiability Checking
  and Symbolic Computation Workshop (SC-Square)}, CEUR Workshop Proceedings
  vol.~3717, pp.~91--110, 2024.

\bibitem{Stoll2025}
T.~Stoll, C.~Qian, B.~Finkelshtein, A.~Parviz, D.~Weber, F.~Frasca,
  H.~Shavit, A.~Siraudin, A.~Mielke, M.~Anastacio, E.~M\"uller,
  M.~Bechler-Speicher, M.~Bronstein, M.~Galkin, H.~Hoos, M.~Niepert,
  B.~Perozzi, J.~T\"onshoff, and C.~Morris.
\newblock {GraphBench}: next-generation graph learning benchmarking.
\newblock arXiv:2512.04475 (v5, May 2026), 2025.

\end{thebibliography}

\appendix

\section{Witnesses and verification}
\label{app:witness}

Both blocks below are preimages of the 20-round SHA-256 compression function
with the standard initial value, feed-forward retained and no padding
constraint.  The verifier is the standard-library script
\path{code/verify_r19.py} from the repository; it shares no code with the
attack.

\begin{footnotesize}
\begin{verbatim}
# all-ones digest (the target of Zaikin, Constraints 2026)
python3 code/verify_r19.py --rounds 20 \
 --hash ffffffffffffffffffffffffffffffffffffffffffffffffffffffffffffffff \
 --words "0b0187bf b9aea692 b66effa5 087dca3a 0caea827 1d2f9916 739a224e a87c0eae
          7f9ef4a7 b318a7de a8848c61 a7a141a4 11ac114b 952299be 97aaa67f c6acfe57"

# solver-generated digest
python3 code/verify_r19.py --rounds 20 \
 --hash d962ca30635f9b74ac6c8c1243a1a9cf800e81bc05f1d2e40764c68c795f7388 \
 --words "a36f4238 f2c9204e b5b6653b 070401f7 5928e4f3 fe766be2 52026907 f7ee1812
          01344603 ea505012 86cbe6cb 6ce73d0b 5c91de6a 43355e3b ff3d5e88 c1ad0b54"
\end{verbatim}
\end{footnotesize}
Each prints \texttt{result: OK}; with \texttt{--rounds 19} or \texttt{--rounds
21} each prints \texttt{result: FAIL}.

\section{Reproducibility}
\label{app:repro}

All code, logs, reviewer artefacts and this paper are in the repository at
\url{https://github.com/kamb-code/sha256-r19-preimage}, under
\texttt{code/} and \texttt{experiments/submask\_2026-09-03/}.

\begin{itemize}
\item \texttt{code/submask\_family.py}: the family, the triangular solve, the
  collapsed test, a table-free algebraic self-test of
  Propositions~\ref{prop:triangular} and~\ref{prop:collapse} (run with no
  arguments), and \texttt{--attack} / \texttt{--measure} modes.
\item \texttt{code/build\_sigma0\_table.py}: builds the $\sig0(u)-u$ table
  (int32, sentinel $-1$, largest root per value; about 215 core-seconds).
\item \texttt{code/gpu\_submask.py}: the PyTorch sweep with up to three root
  tables built on the device; \texttt{code/submask\_triton.py}: the fused
  kernel with \texttt{--verify} (planted preimages) and \texttt{--bench}.
\item \texttt{experiments/submask\_2026-09-03/}: as-run scripts and logs with a
  sha256 manifest; \texttt{FLEET.txt} (hardware, target, seed, exposure,
  and recorded runtime); \texttt{NOTES.md} (corrections found by review);
  \path{R20_PREIMAGE_ALLONES.txt} and \path{R20_PREIMAGE.txt} (witnesses and
  provenance); \path{logs/r19_fixed_targets.log}, \path{logs/frame_search.log} and
  \path{logs/ones_obstruction.log} (the fixed-target runs, the context-family
  search of Table~\ref{tab:frames}, and the all-ones obstruction test);
  \texttt{analysis/} (the corpus mining);
  \texttt{priority/} (the literature search and adjudication);
  \texttt{reviewers/} (separately written automated verification scripts,
  including three SHA-256 implementations written from the specification with
  constants derived from prime roots, and the full reports).
\item Statistical method for nested thresholds: counts $n_k$ of candidates
  whose residual has its low $k$ bits zero are nested.  Under the uniform null,
  $n_{k+1}$ is conditionally Binomial$(n_k,\tfrac12)$.  The manuscript reports
  conditional score summaries and multiplicity-adjusted individual
  transitions; it does not treat the random-denominator ratios as
  unconditionally independent.
\end{itemize}

\end{document}
